% SLPC恒偽式型の証明方法と類題
% 対策プリントの問6（SLPC恒真式型の証明）を踏まえ，逆に恒偽式型（すべての解釈で偽）
% を示すときの方針を整理し，∀・∃が混在する類題3問を解答解説つきでまとめる。
\documentclass[autodetect-engine,dvi=dvipdfmx,ja=standard]{bxjsarticle}

\makeatletter
\providecommand*{\input@path}{}
\g@addto@macro\input@path{{./}{../includes/}}
\makeatother

\input{protocol.tex}
\input{preamble.tex}

\title{SLPC恒偽式型の証明方法と類題}
\author{論理学 問6 関連}
\date{\today}

\begin{document}
\maketitle

対策プリント問6は「\(\forall x[\alpha(x)\supset\beta(x)]\supset(\forall x[\alpha(x)]\supset\forall x[\beta(x)])\)
がSLPC恒真式型（すべての解釈で真）であることを示せ」という問題だった。本資料は，逆に
「SLPC\textbf{恒偽式型}（すべての解釈で偽）であることを示せ」と言われたときの方針を整理し，
\(\forall\)と\(\exists\)が混在する類題を3問，解答解説つきで扱う。

\section{恒真式型と恒偽式型，それぞれの示し方}

\begin{sidebox}[TBteal]{恒真式型の示し方（問6で使った方法）}
任意の解釈 \(I\) を1つ取り，その中で式が真になることを直接追いかける。
\begin{itemize}
  \item \(A\supset B\) を示すときは，「\(A\)が偽なら自動的に真」なので，\(A\)を\textbf{仮定}して\(B\)を示せばよい。
  \item \(\forall x[P(x)]\) を示すときは，領域から任意の1個\(d\)を取り，その\(d\)について\(P(d)\)を示せばよい
    （\(d\)は任意だったので全部について言える）。
  \item \(\forall x[P(x)]\) が真だと分かっているときは，特定の\(d\)について\(P(d)\)を取り出してよい（全称の除去）。
\end{itemize}
\end{sidebox}

\begin{sidebox}[TBred]{恒偽式型の示し方（今回追加する方法）}
恒真式型とは逆に，「すべての解釈で偽になる」ことを示す。2通りのやり方がある。

\textbf{方法A（否定して恒真式型に帰着）：} \(\varphi\)が恒偽式型 \(\Longleftrightarrow\)
\(\sim\varphi\)が恒真式型，という一般的な事実を使い，\(\sim\varphi\)を作って恒真式型の方法（上のボックス）を
そのまま適用する。

\textbf{方法B（背理法・直接示す本筋）：}
\begin{enumerate}
  \item 任意の解釈\(I\)を取り，\(\varphi\)が\(I\)のもとで真だと\textbf{仮定}する。
  \item \(\varphi\)の構造にしたがって仮定を展開する：
    \(\land\)なら両方真，\(\exists x[P(x)]\)が真ならある個体\(d\)で\(P(d)\)が真（witnessを名付ける），
    \(\forall x[P(x)]\)が真ならどの個体についても\(P\)が成り立つ。
  \item これらを組み合わせて，同じ主張について「真」と「偽」を同時に導き，矛盾を得る。
  \item 矛盾が出た以上，1の仮定は誤り。\(\varphi\)は\(I\)のもとで偽。\(I\)は任意だったので恒偽式型。
\end{enumerate}
\end{sidebox}

\begin{sidebox}[TBorange]{\(\sim(A\supset B)\) の展開——ありがちな間違いに注意}
\(A\supset B\;=_{\mathrm{df}}\;\sim A\lor B\) の定義に戻ってからド・モルガンと二重否定の除去を使うと，
\[
  \sim(A\supset B)\;\equiv\;\sim(\sim A\lor B)\;\equiv\;\sim\sim A\land\sim B\;\equiv\;A\land\sim B
\]
となる。\textbf{\(\sim B\)を否定したところが「ならば」ではなく「かつ」になる点に注意}
（\(A\supset\sim B\)や\(\sim A\lor\sim B\)にはならない）。

真理値表で確認すると，\(A\supset B\)が偽になるのは「\(A\)が真で\(B\)が偽」の1行だけなので，
その否定 \(A\land\sim B\) と一致する。

\medskip
\textbf{系：} \(A\supset B\) が恒偽式型であるためには，\(A\land\sim B\) が恒真式型でなければならない。
\(\land\)の恒真性は両方の連言肢が\textbf{別々に}恒真式型であることを要求するので，これは
\textbf{「\(A\)が恒真式型」かつ「\(B\)が恒偽式型」}という2つの独立した事実に分解できる
（\(T\supset F\)だけが\(F\)になる組合せなので，\(A\supset B\)が常に偽であるためには
\(A\)が常に真，\(B\)が常に偽である必要がある，という言い換えでもある）。
\end{sidebox}

\section{参考：問6の解答（恒真式型の型を確認する）}

\begin{exercise}{問6（参考）}{ref-q6}
次のSLPC論理式型がSLPC恒真式型であることを示せ。
\[
  \forall x[\alpha(x)\supset\beta(x)]\supset\bigl(\forall x[\alpha(x)]\supset\forall x[\beta(x)]\bigr)
\]
\end{exercise}
\begin{solution}{問6（参考）}{ref-sol-q6}
任意の解釈\(I\)を取る。前件\(\forall x[\alpha(x)\supset\beta(x)]\)が偽なら全体は自動的に真。
前件が真であると仮定し，後件も真であることを示す。後件も\(\supset\)の形なので，
さらに\(\forall x[\alpha(x)]\)が真であると仮定し，\(\forall x[\beta(x)]\)が真であることを示せばよい。

領域\(D\)の任意の元\(d\)を取る。仮定より\(\alpha(d)\supset\beta(d)\)と\(\alpha(d)\)がともに真だから，
分離規則（modus ponens）で\(\beta(d)\)も真。\(d\)は任意だったので\(\forall x[\beta(x)]\)が真。
以上より式全体は\(I\)のもとで常に真，\(I\)は任意だったのでSLPC恒真式型である。\hfill\(\blacksquare\)
\end{solution}

\section{類題1：\(\forall\)と\(\exists\)がそのまま矛盾する基本形}

\begin{exercise}{類題1}{ex1}
次の式がSLPC恒偽式型であることを示せ。
\[
  \forall x[\alpha(x)]\;\land\;\exists x[\sim\alpha(x)]
\]
\end{exercise}

\begin{solution}{類題1}{sol1}
\textbf{方法B（背理法）：}

任意の解釈\(I\)を取り，上式が\(I\)のもとで真だと仮定する。\(\land\)なので
\(\forall x[\alpha(x)]\)と\(\exists x[\sim\alpha(x)]\)がともに真である。

\begin{itemize}
  \item \(\exists x[\sim\alpha(x)]\)が真だから，ある個体\(d\in D\)で\(\sim\alpha(d)\)が真である
    （存在の除去：この\(d\)を名付けて以後使ってよい）。
  \item \(\forall x[\alpha(x)]\)が真だから，\textbf{その同じ}\(d\)についても\(\alpha(d)\)が真である
    （全称の除去）。
\end{itemize}
\(\alpha(d)\)と\(\sim\alpha(d)\)が同時に真というのは矛盾である。したがって最初の仮定は誤りであり，
上式は\(I\)のもとで偽である。\(I\)は任意の解釈だったから，この式はすべての解釈のもとで偽，
すなわちSLPC恒偽式型である。\hfill\(\blacksquare\)

\textbf{方法A（否定して恒真式型を示す・別解）：}

\(\sim\bigl(\forall x[\alpha(x)]\land\exists x[\sim\alpha(x)]\bigr)\)をド・モルガンで展開する。
\[
  \sim\bigl(\forall x[\alpha(x)]\land\exists x[\sim\alpha(x)]\bigr)
  \;\equiv\;
  \sim\forall x[\alpha(x)]\;\lor\;\sim\exists x[\sim\alpha(x)]
\]
ここで量化子の否定則 \(\sim\forall x[P(x)]\equiv\exists x[\sim P(x)]\)，
\(\sim\exists x[P(x)]\equiv\forall x[\sim P(x)]\) を使うと，
\[
  \equiv\;
  \exists x[\sim\alpha(x)]\;\lor\;\forall x[\sim\sim\alpha(x)]
  \;\equiv\;
  \exists x[\sim\alpha(x)]\;\lor\;\forall x[\alpha(x)]
\]
ここで\(\forall x[\alpha(x)]\)を1つの命題\(P\)とみなせば，この式は\(\sim P\lor P\)，
すなわちPM-Th.3（排中律）をSLPCの命題\(\forall x[\alpha(x)]\)に持ち上げただけの形であり，恒真式型である。
よって元の式はその否定なので恒偽式型である。\hfill\(\blacksquare\)
\end{solution}

\section{類題2：問6と同じ「全称の除去→分離規則」を使う版}

\begin{exercise}{類題2}{ex2}
次の式がSLPC恒偽式型であることを示せ。
\[
  \forall x[\alpha(x)]\;\land\;\exists x[\beta(x)]\;\land\;\forall x[\alpha(x)\supset\sim\beta(x)]
\]
\end{exercise}

\begin{solution}{類題2}{sol2}
任意の解釈\(I\)を取り，上式が\(I\)のもとで真だと仮定する。\(\land\)なので3つの連言肢
\(\forall x[\alpha(x)]\)，\(\exists x[\beta(x)]\)，\(\forall x[\alpha(x)\supset\sim\beta(x)]\)が
すべて真である。

\begin{itemize}
  \item \(\exists x[\beta(x)]\)が真だから，ある個体\(d\in D\)で\(\beta(d)\)が真である（witnessの導入）。
  \item \(\forall x[\alpha(x)]\)が真だから，その\(d\)について\(\alpha(d)\)も真である（全称の除去）。
  \item \(\forall x[\alpha(x)\supset\sim\beta(x)]\)が真だから，その\(d\)について
    \(\alpha(d)\supset\sim\beta(d)\)も真である（全称の除去）。
  \item \(\alpha(d)\)と\(\alpha(d)\supset\sim\beta(d)\)がともに真なので，分離規則（modus ponens）より
    \(\sim\beta(d)\)も真である。
\end{itemize}
しかし最初にwitnessとして取った\(d\)は\(\beta(d)\)が真であった。\(\beta(d)\)と\(\sim\beta(d)\)が
同時に真というのは矛盾である。したがって最初の仮定は誤りであり，上式は\(I\)のもとで偽である。
\(I\)は任意だったから，SLPC恒偽式型である。\hfill\(\blacksquare\)

（問6の証明で使った「\(\forall\)を外して分離規則で\(\beta(d)\)を導く」操作と全く同じ部品を使い，
最後に矛盾へ着地させただけの構造になっている。）
\end{solution}

\section{類題3：\(A\supset B\)全体が恒偽式型になる形}

\begin{exercise}{類題3}{ex3}
次の式がSLPC恒偽式型であることを示せ。
\[
  \forall x[\alpha(x)\lor\sim\alpha(x)]
  \;\supset\;
  \bigl(\forall x[\alpha(x)]\land\exists x[\sim\alpha(x)]\bigr)
\]
\end{exercise}

\begin{solution}{類題3}{sol3}
全体は\(A\supset B\)の形（\(A=\forall x[\alpha(x)\lor\sim\alpha(x)]\)，
\(B=\forall x[\alpha(x)]\land\exists x[\sim\alpha(x)]\)）である。
\(A\supset B\)が恒偽式型であるためには，\(A\)が恒真式型かつ\(B\)が恒偽式型であればよい
（\(T\supset F\)だけが偽になる組合せであり，これがすべての解釈で成り立つには
\(A\)が常に真，\(B\)が常に偽である必要がある）。

\textbf{①\(A\)が恒真式型であることを示す：}
任意の解釈\(I\)と領域\(D\)の任意の元\(d\)について，\(\alpha(d)\lor\sim\alpha(d)\)は
命題論理の排中律（PM-Th.3）の代入例として真である。\(d\)は任意だったので
\(\forall x[\alpha(x)\lor\sim\alpha(x)]\)は\(I\)のもとで真。\(I\)は任意だったので恒真式型である。

\textbf{②\(B\)が恒偽式型であることを示す：}
\(B=\forall x[\alpha(x)]\land\exists x[\sim\alpha(x)]\)は類題1そのものであり，
すでに恒偽式型であることを示した。

\textbf{③組み合わせる：}
①・②より，任意の解釈\(I\)のもとで\(A\)は真，\(B\)は偽である。したがって
\(A\supset B\)は\(I\)のもとで\(\true\supset\false=\false\)。\(I\)は任意だったから，
元の式はSLPC恒偽式型である。\hfill\(\blacksquare\)

（\(A\land\sim B\)が恒真式型であることを1発で示すよりも，\(A\)の恒真性と\(B\)の恒偽性を
別々に示して組み合わせるほうが見通しがよく，実際に類題1の結果がそのまま部品として使えている。）
\end{solution}

\end{document}
